\scanpage{9}\mspage{7}
On se donne un ens. $M_0$, et
\[
\left\{\begin{aligned}
M_0\times M_0&\overset{\mathrm{Ch}}{\longrightarrow}\unsure{\mathscr H|\mathcal H}&\text{noté }\quad(a,b)&\longrightarrow\mathrm{Ch}(a,b)\\
&&&\text{ou }\mathrm{Hom}\,\mathrm{ext}_{\mathscr H}(a,b)\\
M_0\times M_0\times M_0&\overset{\mathrm{comp}}{\longrightarrow}\mathrm{Fl}\,\mathscr H &(a,b,c)&\mapsto\bigl(\mathrm{Ch}(a,b)\times\mathrm{Ch}(b,c)\\
&&&\hspace{3em}\longrightarrow\mathrm{Ch}(a,c)\bigr)
\end{aligned}\right.
\]
avec deux axiomes

a) Associativité pour les compositions, pour $(a,b,c,d)\in M_0^{\unsure{4}}$
\[
\begin{aligned}
&\deleted{\unsure{M_0\times M_0\times M_0\times M_0}\longrightarrow}\\
&\deleted{\mathrm{Ch}(a,b)\times\mathrm{Ch}(b,c)\times\mathrm{Ch}(c,d)\longrightarrow\mathrm{Ch}(a,d)}
\end{aligned}
\]
\diagram{\begin{tikzcd}[ampersand replacement=\&,column sep=small]
\mathrm{Ch}(a,b)\times\mathrm{Ch}(b,c)\times\mathrm{Ch}(c,d) \arrow[r] \arrow[d] \& \mathrm{Ch}(a,c)\times\mathrm{Ch}(c,d) \arrow[d]\\
\mathrm{Ch}(a,b)\times\mathrm{Ch}(b,d) \arrow[r] \& \mathrm{Ch}(a,d)
\end{tikzcd}}

b) Existence d'objets unité bilatères
\[
1_a\in\pi_0(\mathrm{Ch}(a,a))\qquad\text{pour tt }a\in M_0.
\]
Alors les $\pi_0(\mathrm{Ch}(a,b))$ et leur composition définissent sur $M_0$\deleted{\unsure{,}} \inserted{comme ens. d'objets} une structure de \uline{catégorie} \inserted{$M_{0\pi}$}, ayant les $1_a$ comme identités, et on trouve que
\[
(a,b)\unsure{\vdash|\mapsto}\mathrm{Ch}(a,b)
\]
peut s'interpréter \unsure{à partir|: partir} d'un bi\uline{foncteur} canonique
\[
M_{0\pi}^{o}\times M_{0\pi}\longrightarrow\mathscr H
\]
dont le composé avec
\[
\mathscr H\xrightarrow{\pi_0}\mathrm{Ens}
\]
n'est autre que le bifoncteur $(a,b)\longrightarrow\mathrm{Hom}_{M_{0\pi}}(a,b)$,
% archive sheet number: 8

